\section{All higher symplectic cores in even characteristic}
\label{sec:hes-all-even}

\paragraph{Status and scope.}
This is the complete author derivation D-HES-04--06, with root and
independent-peer whole-proof review both passing. The reviewed
mathematical text has SHA-256
\texttt{ab4898d2725c48b0993539ead8e56c27b9630f0433aee185514c48ab1546ae18}.
The author exact certificate passes, and the peer independently
reconstructed the finite-field identities with different arithmetic.
It extends the reviewed D-HES-01--03
from $\mathbf F_2$ to every even field. It treats all ordinary irreducible
modules and all compatible linear normalizers, including field operators,
scalars and permutations of tensor factors. The conclusion concerns the
actual nine-dimensional period image in the sixfold-source problem.
Neither a reflection hypothesis nor a determinant hypothesis is imposed.

\subsection*{D-HES-04. Core support over every even field}

Let $q\ge4$ be a power of two, $n\ge4$, and
$S=\operatorname{Sp}_{2n}(q)$. It is centerless simple. Its Schur multiplier
is trivial in this range: the ordinary symplectic multiplier is
$(2,q-1)=1$, and there is no exception here in the complete list
\cite[Chapter~5, Table~1, numbered p.~35]{BI:HHM}.
For a finite-order operator $x$ on a complex space of dimension $a$, put
\[
              \nu(x)=a-\max_{\lambda}\dim\ker(x-\lambda I).
\]
This quantity is invariant under multiplying $x$ by a scalar.

We use the following precise published unipotent estimate:
\[
 |\chi(u)|\le c(q)\chi(1),\qquad
 c(4)=\frac34,\qquad
 c(q)=\frac1{\sqrt q-1}\quad(q\ge8),
 \qquad 1\ne u\text{ unipotent},\quad1\ne\chi\in\operatorname{Irr}(S).
\tag{HES.12}
\]
This is \cite[Theorem~2.4(i), definitions and statement on author
pp.~2--4]{HES:Liebeck2017}. The subtraction of $1$ is outside the
radical. In particular $c(8)$ is not less than $1/2$. The uniform
comparison needed below is instead
\[
              c(q)\le\frac1{\sqrt8-1}<\frac9{16}
                  \quad(q\ge8),
       \qquad 625<648.
\tag{HES.13}
\]
The last integer inequality is equivalent to $25/9<\sqrt8$.

Every nontrivial irreducible $S$-module is faithful, and has degree
\[
                          a\ge q^{n-1}-1.
\tag{HES.14}
\]
Indeed embed $\operatorname{GL}_n(q)$ by
$A\mapsto\operatorname{diag}(A,A^{-\mathsf T})$, and in it take the row
subgroup
\[
 U_0=\left\{\begin{pmatrix}1&v\\0&I_{n-1}\end{pmatrix}:
                   v\in\mathbf F_q^{n-1}\right\}.
\]
Its Levi subgroup $\operatorname{GL}_{n-1}(q)$ is transitive on all
nontrivial complex additive characters of $U_0$. Faithfulness ensures
that one, hence all $q^{n-1}-1$, occur with the same positive multiplicity.
For $q=4$ we also have $a\ge378$, by restricting to the block subgroup
$\operatorname{Sp}_6(4)$: its minimum nontrivial degree $378$ is the
complete exact character-table assertion certified in D-ESC-01.
Restriction has a nontrivial constituent because the original module
is faithful.

\begin{lemma}[D-HES-04a]
Let $h\in S$ have odd prime order $p$ and stabilize a nonzero totally
isotropic subspace of its natural module. There is a nonzero
$h$-normalized elementary abelian $2$-subgroup $U\le S$ on which $h$
has no nonidentity fixed point, and $|U|\ge4$.
\end{lemma}

\begin{proof}
The construction of D-HES-02a works over $\mathbf F_q$, with its
underlying additive groups viewed as $\mathbf F_2$-spaces. We give the
details needed to avoid a field-linearity assumption on the resulting
$U$. Write $b$ for the natural symplectic form and $R$ for an invariant
isotropic space with basis $r_1,\ldots,r_s$.
If $A=h|_R\ne I$, define for every symmetric $B=(b_{ij})$ over
$\mathbf F_q$
\[
               N_B(x)=\sum_{i,j}b_{ij}b(r_j,x)r_i.
\]
Then $N_BN_C=0$, and $1+N_B$ is symplectic. These elements form an
elementary abelian $2$-group; conjugation by $h$ sends $B$ to
$ABA^{\mathsf T}$. This action is nontrivial, since
$B=vv^{\mathsf T}$ goes to $(Av)(Av)^{\mathsf T}$ and the diagonal
entries $v_i^2$ recover $v$ uniquely in a finite field of characteristic
two. Thus $A\ne I$ moves a rank-one matrix.

If instead $h$ fixes $R$ pointwise, choose $0\ne e\in R$ and $f_0$
with $b(e,f_0)=1$. The vector $f=\sum_{j=0}^{p-1}h^jf_0$ is fixed and
$b(e,f)=1$. Hence the natural module is the $h$-stable orthogonal sum
$\langle e,f\rangle\perp W_0$, with nontrivial action on $W_0$.
For $v\in W_0,c\in\mathbf F_q$ define
\[
 q(v,c)e=e,\qquad q(v,c)w=w+b(w,v)e,\qquad q(v,c)f=f+v+ce.
\]
These are symplectic and obey
\[
          q(v,c)q(w,d)=q(v+w,c+d+b(v,w)).
\tag{HES.15}
\]
Because $b$ is alternating and symmetric in characteristic two, this
is an abelian group of exponent two, of order $q^{2n-1}$. Conjugation
by $h$ is $q(v,c)\mapsto q(hv,c)$, a nontrivial action.

In either case the elementary abelian group is a semisimple
$\mathbf F_2\langle h\rangle$-module, since $p$ is odd. Take its sum of
nontrivial simple summands. This gives a nonzero invariant subgroup
$U$ with $C_U(h)=1$. It has dimension at least two over $\mathbf F_2$,
because $\operatorname{GL}_1(2)=1$. No assertion that $U$ is an
$\mathbf F_q$-subspace is necessary.
\end{proof}

\begin{lemma}[D-HES-04b]
For every nontrivial irreducible complex $S$-module of degree $a$,
\[
 \nu(\rho(x))\ge a/8\quad(q=4),\qquad
 \nu(\rho(x))>7a/32\quad(q\ge8),\qquad 1\ne x\in S.
\tag{HES.16}
\]
\end{lemma}

\begin{proof}
First take odd-prime-order $h$ stabilizing an isotropic space, and use
$U$ from D-HES-04a. All its nonidentity elements are unipotent. If
$t=\dim V^U$, averaging (HES.12) yields
\[
       t\le a\left(c(q)+\frac{1-c(q)}{|U|}\right).
\]
The dual $h$-action on the character group of $U$ has no nontrivial
fixed character; its nontrivial orbits all have length $p$. On an
orbit-sum of weight spaces the eigenvalues of $h$ occur uniformly,
because $h^p=1$. A specified eigenvalue therefore has multiplicity at
most $t+(a-t)/p$, giving
\[
 \nu(\rho(h))\ge(1-1/p)(a-t)
       \ge\frac23\frac34(1-c(q))a=\frac{1-c(q)}2a.
\tag{HES.17}
\]
For an involution, its two eigenvalue multiplicities give exactly
$\nu(\rho(h))=(a-|\chi(h)|)/2$, with the same lower bound.

It remains to treat an odd-prime-order $h$ which stabilizes no nonzero
totally isotropic subspace. Its primary decomposition has no pair of
distinct reciprocal irreducible polynomials: either member would give
an invariant totally isotropic space. Its $1$-eigenspace is zero,
since a nonzero such space is nondegenerate symplectic and contains
an invariant isotropic line. Every remaining primary polynomial is
self-reciprocal of degree $2d$.

On its primary summand set $E=\mathbf F_{q^{2d}}$ and
$E_0=\mathbf F_{q^d}$. The adjoint involution on $E$ is
$z\mapsto\bar z=z^{q^d}$. The nondegenerate trace pairing uniquely
defines $H$ by
\[
 b(av,w)=\operatorname{Tr}_{E/\mathbf F_q}(aH(v,w))
                       \quad(a\in E).
\]
Adjointness and symmetry give
$H(av,bw)=a\bar bH(v,w)$ and
$H(w,v)=\overline{H(v,w)}$, so this is a nondegenerate Hermitian form.
If the $E$-dimension were at least two, diagonalization and surjectivity
of $E^\times\to E_0^\times$ by the norm would produce an isotropic
vector. Its entire $E$-line is totally isotropic for $b$, a contradiction.
Every primary multiplicity is thus one. The centralizer is a product
of norm-one tori, so
\[
               |C_S(h)|=\prod_i(q^{d_i}+1)\le(q+1)^n,
                      \qquad\sum_i d_i=n.
\tag{HES.18}
\]
Column orthogonality and (HES.14) give
\[
          \frac{|\chi(h)|}{a}
             \le\frac{(q+1)^{n/2}}{q^{n-1}-1}<\frac12.
\tag{HES.19}
\]
To check this uniformly, at $n=4,q\ge4$ the right side is at most
$100/(63q)\le25/63<1/2$: use
$(q+1)^2\le25q^2/16$ and $q^3-1\ge63q^3/64$.
As $n$ increases its ratio decreases, since
$(q^n-1)/(q^{n-1}-1)>q>\sqrt{q+1}$.
Every nonidentity power of $h$ is again elliptic. Fourier inversion
therefore bounds every eigenvalue multiplicity by
$a(1/2+1/(2p))\le2a/3$, giving the stronger $\nu(\rho(h))\ge a/3$.

For an arbitrary nonidentity $x$, choose a nonidentity prime-order
power $h=x^k$. Each $x$-eigenspace lies in an $h$-eigenspace, so
$\nu(\rho(x))\ge\nu(\rho(h))$. Equations (HES.12)--(HES.13) and
(HES.17)--(HES.19) prove the assertion.
\end{proof}

\subsection*{D-HES-05. All compatible normalizer operators}

\begin{lemma}[D-HES-05a]
For $q\ge8$, every nonscalar finite-order linear operator $T$ normalizing
an irreducible $S$-module of degree $a$ has $\nu(T)>7a/64$.
\end{lemma}

\begin{proof}
The induced automorphism $\alpha$ of $S$ is nontrivial, since the
centralizer of an irreducible representation consists of scalars.
Choose $s\in S$ with $\alpha(s)\ne s$. Then
$[T,\rho(s)]=\rho(\alpha(s)s^{-1})$ is a nonidentity core operator;
faithfulness and centerlessness ensure it is nonscalar. For an
eigenvalue $\lambda$ of $T$ of maximum multiplicity,
\[
 \begin{split}
 \nu([T,\rho(s)])
 &\le\operatorname{rank}([T,\rho(s)]-I)\\
 &=\operatorname{rank}(T\rho(s)-\rho(s)T)\\
 &\le2\operatorname{rank}(T-\lambda I)=2\nu(T).
 \end{split}
\]
Apply D-HES-04b. This proof covers every automorphism, without a
restriction on the order or type of the outer automorphism.
\end{proof}

\begin{lemma}[D-HES-05b: the field-involution endpoint]
For $S=\operatorname{Sp}_{2n}(4)$, $n\ge4$, every nonscalar finite-order
linear normalizer operator on an irreducible degree-$a$ module satisfies
\[
                              \nu(T)\ge7a/64.
\tag{HES.20}
\]
\end{lemma}

\begin{proof}
The outer automorphism group is $C_2$, by
\cite[Section~2.2.2, Table~2.1, author p.~14]{HES:LiXia2022}.
The entrywise field automorphism $\sigma:z\mapsto z^2$ represents its
nontrivial element. To check the latter without a naming convention,
embed the determinant-one matrix
$\left(\begin{smallmatrix}\omega&1\\1&0\end{smallmatrix}\right)$
in a symplectic two-plane, where $\omega^2+\omega=1$.
Its natural trace changes from $\omega$ to $\omega^2$ under $\sigma$,
so $\sigma$ is not inner. Thus every automorphism is
$\operatorname{inn}(g)$ or $\operatorname{inn}(g)\sigma$.

In the inner case, $T$ is a scalar times $\rho(g)$ and D-HES-04b gives
$\nu(T)\ge a/8$. In the outer case put
$\alpha=\operatorname{inn}(g)\sigma$. If $g\sigma(g)\ne1$, then
$T^2$ is a scalar times the nonidentity core operator
$\rho(g\sigma(g))$, and
$\nu(T)\ge\nu(T^2)\ge a/8$. We may therefore suppose
$g\sigma(g)=1$. The natural semilinear operator $s=g\sigma$ then
satisfies $s^2=1$.

Here is an explicit symplectic descent for this last case. Write
$W_*=\operatorname{Fix}(s)$. For any natural vector $v$, both
\[
             a=v+s(v),\qquad b=\omega^2v+\omega s(v)
\]
belong to $W_*$, and $v=b+\omega a$. If $u+\omega v=0$ with
$u,v\in W_*$, applying $s$ and subtracting gives $v=0$, then $u=0$.
Consequently the underlying $\mathbf F_2$-space is
$W_*\oplus\omega W_*$, with $\dim_{\mathbf F_2}W_*=2n$; an
$\mathbf F_2$-basis of $W_*$ is an $\mathbf F_4$-basis of the natural
module. Its symplectic pairing is $\mathbf F_2$-valued, since
$b(sx,sy)=b(x,y)^2$, and is nondegenerate since $W_*$ spans over
$\mathbf F_4$. A symplectic $\mathbf F_2$-basis therefore conjugates
$s$ to the entrywise $\sigma$ by an element of $S$.
After this conjugation the operator $T$ implements $\sigma$ on $\rho(S)$.
Its square is scalar, so rescale it to $T^2=1$.

Use $U_0\cong\mathbf F_4^r$, $r=n-1\ge3$, from (HES.14).
Its characters are
\[
        \chi_a(v)=(-1)^{\operatorname{Tr}_{\mathbf F_4/\mathbf F_2}(a\cdot v)}
                         \quad(a\in\mathbf F_4^r).
\]
All nontrivial characters have the same multiplicity $m$, by the
$\operatorname{GL}_r(4)$ action. Put $N=4^r$ and let $t$ be the
trivial weight multiplicity, so $a=t+(N-1)m$. The unipotent bound gives
\[
 t\le a\left(\frac34+\frac1{4N}\right),
                 \qquad m\ge\frac{a}{4N}.
\tag{HES.21}
\]
On the dual group, $\sigma$ sends the parameter to its coordinatewise
square. Exactly $M=2^r$ parameters are fixed. The other $N-M$ weights
form pairs. On each paired sum, $T^2=1$ forces equal $+$ and $-$
multiplicities, both equal to $m$. Thus either eigenvalue of $T$ misses
at least $(N-M)m/2$ dimensions, and
\[
       \nu(T)\ge\frac{N-M}{2}m
            \ge\frac{1-2^{-r}}8a\ge\frac7{64}a.
\tag{HES.22}
\]
The argument bounds actual eigenspaces of this field operator; it does
not assume that any special eigenvalue bound survives an arbitrary
scalar convolution.
\end{proof}

\subsection*{D-HES-06. The complete even-characteristic higher-rank exclusion}

\begin{theorem}[D-HES-06]
An irreducible inducing image in the signature-$(3,3)$ sixfold-source
problem cannot have a nine-dimensional actual period image if it has
an irreducible normal subgroup whose quasisimple factors have quotients
$\operatorname{Sp}_{2n_i}(q_i)$ with $n_i\ge4$ and all $q_i$ even.
The $q_i$ may differ. Arbitrary compatible scalar extensions, outer
operators, factor permutations and larger inducing character fields
are allowed.
\end{theorem}

\begin{proof}
The trivial multipliers in this range identify the core with a direct
product of these simple groups. Write its irreducible module as
$V=\bigotimes_i V_i$, with nontrivial $V_i$ of degree $a_i$ and
$e=\prod_i a_i$. A normalizer permutes simple factors and is, up to
scalar, a tensor of intertwining maps followed by this permutation.
For a factor-preserving operator the normalized eigenvalue
multiplicities convolve; their maximum cannot exceed the maximum for
any nonscalar factor. D-HES-02--03 and D-HES-05 give the common support
bound $7e/64$.

For completeness, a nontrivial cycle of length $\ell\ge2$ on equal
degree-$a$ factors can be conjugated so that all twists occur on one
coordinate. Their product is diagonalizable: a power of the whole
operator is a tensor of these products, and finite order excludes a
nontrivial unipotent Jordan part. In an eigenbasis this is a weighted
cyclic permutation of $a^\ell$ words. There are $a$ constant words;
every other word orbit has length at least two, contributing at most
one copy of a specified eigenvalue. The maximum multiplicity is at
most $a+(a^\ell-a)/2\le3a^\ell/4$. Tensoring other cycles cannot
increase the normalized maximum. Such an operator has support at
least $e/4$, so again at least $7e/64$. Nonidentity scalars have actual
moving rank $e$.

By Gow's theorem \cite[p.~583 and Section~3]{HES:Gow1981}, all the
factor characters are real. Hence $V$ is self-dual and has a unique
line of invariant nondegenerate bilinear forms. The inducing group
$J$ acts on that line through a multiplier $\mu$. It is nontrivial:
otherwise all finite-order $J$-operators would have reciprocal
eigenvalue pairs and real trace, whereas the original imaginary
quadratic field is contained in the inducing character field. This
uses the proven field inclusion from D-BI-01, without assuming that
the two fields coincide.

The actual intermediate-cover variation has Euler formula
\[
                   6=2(h-1)e+\sum_{j=1}^r\operatorname{rank}(g_j-I).
\tag{HES.23}
\]
Every factor has degree at least $35$: use the verified $q=2$ endpoint,
the $q=4$ bound $378$, or (HES.14) for $q\ge8$.
As in D-HES-03, $h\ge2$ is impossible and $h=1$ has actual image
dimension at most six. A nine-dimensional image therefore requires
$h=0$, $r\ge12$ and total moving rank $2e+6$.
In a generating product-one tuple at least two inertias have
nontrivial multiplier. Their $1$-eigenspaces are totally isotropic
and they each have moving rank at least $e/2$. The other at least
ten have moving rank at least $7e/64$. Consequently
\[
               2e+6\ge e+10\frac7{64}e=\frac{67}{32}e,
                            \qquad e\le64.
\tag{HES.24}
\]
Two or more factors would give $e\ge35^2>64$. For a single factor,
$q=4$ gives $e\ge378$, and $q\ge8$ gives
$e\ge q^{n-1}-1\ge511$, again impossible. The remaining single
$q=2$ case is already excluded by D-HES-03, whose sharper support
bound forces $e\le24<35$. This exhausts every case.
\end{proof}

\paragraph{Verification boundary.}
The new universal arguments are the elementary $2$-subgroup construction,
the elliptic trace estimate, the normalizer commutator estimate and the
explicit $\mathbf F_4/\mathbf F_2$ symplectic descent with weight pairs.
Their proofs are written above. A separate exact finite-field program
checks the displayed identities and endpoint inequalities; it does not
replace the universal proofs. The classification and unipotent inputs
are cited with exact scopes. No unavailable Gluck--Magaard (1998)
statement is used, and the computed $\operatorname{Sp}_{12}(2)$ table
remains redundant.
